% conclusion.tex
%
% Phase 5 synthesis (v1.0).  The conclusion is structured as:
%   - subsec:results_summary         -- what each phase established
%   - subsec:terminal_characterisation -- the framework's terminal prediction
%   - subsec:quantitative_predictions  -- the testable numbers
%   - subsec:methodological           -- substrate-first, debt to measurement
%   - subsec:open_directions          -- follow-on papers and open questions

\label{sec:conclusion}

\subsection{What Phase~1 through Phase~4 established}
\label{subsec:results_summary}

Eq.~\eqref{eq:automorphism_conjecture} of Paper~I is settled to
\texttt{PART} (characterisation outcome).  The four phases of the
work plan (\nolinkurl{notes/work_plan.md}) close as follows.

\paragraph{Phase~1 -- Chirality alignment (\texttt{PART}).}  Under
Branch~A SU(3) and the existing bipartite tick rule, the unique
linear bipartite parity is $P = \sigma_x \otimes I_2 \otimes I_3$.
It anticommutes with $\gamma_5 = \sigma_z \otimes I_2 \otimes I_3$,
swaps the SM's $\psi_L \leftrightarrow \psi_R$, and commutes with
every proposed $SU(2)_W$ generator.  The lattice's bipartite
$\mathbb{Z}_2$ and the SM's chirality $\mathbb{Z}_2$ are therefore
orthogonal Bloch involutions on the chirality $\mathbb{C}^2$:
spatial parity, not chirality projection.  Symbolically verified
by \nolinkurl{src/utilities/chirality_parity_alignment.py}.

\paragraph{Phase~1.5 -- Modified tick rule for SM $CP$ (\texttt{FAIL}).}
The natural antilinear modification
$T_\text{ext}^B = is\,I_{12} + c\,(\sigma_x \otimes I_2 \otimes C) \circ K$
that would make Branch~B SU(3) ($\mathbf{3} \oplus \bar{\mathbf{3}}$) a
global symmetry requires the colour matrix $C$ to anticommute with
every Gell-Mann generator.  No non-zero $C$ does: anticommutation
with $\lambda_3$ leaves only $C_{12}, C_{21}, C_{33}$ unconstrained,
and anticommutation with $\lambda_8$ kills those.  The lattice
cannot incorporate SM-style $CP$ as a discrete symmetry of any
natural tick modification.  Symbolically verified by
\nolinkurl{src/utilities/tick_rule_cp_modified.py}; the
representation-theoretic origin is that $\mathbf{3}$ and
$\bar{\mathbf{3}}$ are inequivalent irreps of SU(3).

\paragraph{Phase~2 -- Non-abelian $SU(3)$ from $\mathbb{C}^3$ memory (\texttt{PASS}).}
The natural rotate-toward-$|j\rangle$ colour-memory tick rule
($X_1 = \lambda_1 + \lambda_4$ for $j = 1$, with $S_3$ conjugates)
closes under Lie-bracket closure to the full $\mathfrak{su}(3)$.
Restricted candidates (diagonal-only, single-off-diagonal) yield
strict subalgebras (Cartan $\mathfrak{h}$ or one $\mathfrak{su}(2)$
subalgebra).  The framework's $\mathfrak{su}(3)$ factor is
\emph{dynamically generated} by the lattice for any natural
real-symmetric colour-memory tick rule; the discrete RGB
$\mathbb{Z}_3$ of \nolinkurl{automorphism_rgb_su3.py} is recovered as
the abelian-Cartan corner of this larger structure.  Symbolically
verified by \nolinkurl{src/utilities/su3_generation_from_colour_memory.py}.

\paragraph{Phase~3 -- Containment vs equality on $\mathrm{Aut}_\text{ext}$ (\texttt{PART}).}
The discrete-Hermitian centralizer of the bipartite tick rule on
$\mathbb{C}^{12}$ is
\begin{equation}
\label{eq:su6_su6_u1}
\mathfrak{Aut}_{\text{centralizer}} \;=\; \mathfrak{su}(6)_+ \oplus \mathfrak{su}(6)_- \oplus \mathfrak{u}(1),
\qquad \dim = 35 + 35 + 1 = 71,
\end{equation}
with the two $\mathfrak{su}(6)$ factors acting on the $\pm 1$
chirality eigenspaces and $J_1$ as the central $\mathfrak{u}(1)$.
The SM gauge subalgebra $\mathfrak{su}(3) \oplus \mathfrak{su}(2)
\oplus \langle J_1 \rangle$ (dim 12 in $\mathfrak{su}(12)$) embeds
diagonally via the $\mathbf{6} = (\mathbf{3}, \mathbf{2})$
branching; the 59 ``extras'' decompose under $SU(3) \times SU(2)$
as
\begin{equation}
\label{eq:extras_decomp}
\text{59 extras} \;=\; 48 \times (\mathbf{8}, \mathbf{3})_{\text{leptoquark}}
\;+\; 11 \times \text{chirality-}\sigma_x\text{ shadow SM}.
\end{equation}
Bracket-structure verification
(\nolinkurl{aut_centralizer_extras_commutators.py}):
$[\text{extras}, \text{SM}] \subseteq \text{extras}$ (0 of 708
brackets land in SM); $[\text{extras}, \text{extras}]$ has 226 of
1711 brackets in SM, so the extras are an SM-module but not a Lie
ideal --- exactly the algebraic shape of a left-right symmetric
broken gauge symmetry pattern.  Symbolically verified by
\nolinkurl{src/utilities/automorphism_centralizer_extended.py}.

\paragraph{Phase~4 -- Wilson $1/g^2$ for $SU(2)_W$, $SU(3)$ (\texttt{PART}).}
Paper~I's bipartite-plaquette Q-tensor (eigenvalues
$\{4, 4, 16\}$) is inherited unchanged for non-abelian gauge
groups; the SU(N) trace identity introduces only a global
$T_F = 1/2$ factor.  Spectator-factor counting on the per-site
$\mathbb{C}^{12}$ gives the framework's quantitative gauge-coupling
prediction:
\begin{equation}
\label{eq:lattice_scale_ratios}
\frac{1}{g_1^2} : \frac{1}{g_2^2} : \frac{1}{g_3^2} \;=\; 12 : 3 : 2 \quad \text{at the lattice scale } 1/a,
\end{equation}
equivalently $g_1^2 : g_2^2 : g_3^2 = 1 : 4 : 6$, and the sharp
ratio
\begin{equation}
\label{eq:g3_over_g2}
g_3^2 / g_2^2 = 3/2 \quad \text{at the lattice scale,}
\end{equation}
independent of the universal one-loop prefactor $c$ (which
Paper~I and Paper~II both leave open as the explicit
$-\mathrm{Tr}\,\ln D_\text{lat}[U]$ calculation).  Symbolically
verified by \nolinkurl{src/utilities/induced_gauge_action_nonabelian.py}.

\subsection{The framework's terminal characterisation}
\label{subsec:terminal_characterisation}

The conjunction of Phase~1 through Phase~4 settles the conjecture's
central row to \texttt{PART} with a precise structural reading.
Eq.~\eqref{eq:automorphism_conjecture} holds as containment
\begin{equation}
\label{eq:containment_again}
\mathrm{Aut}(\mathcal{T}_\diamond^3, \mathcal{A}=1)
\;\supseteq\;
SO(3,1) \times SU(3) \times SU(2) \times U(1),
\end{equation}
where the right-hand side is the \emph{factor-product-irreducible
subalgebra} of the lattice's discrete-Hermitian automorphism
centralizer of Eq.~\eqref{eq:su6_su6_u1}.  Equality $=$ does
\emph{not} hold at the discrete level: the centralizer is 71-dim,
not 18-dim; the SM gauge group is exactly the projection that
respects the tensor-factor decomposition
$\mathbb{C}^{12} = \mathbb{C}^2_\text{chir} \otimes
\mathbb{C}^2_\text{iso} \otimes \mathbb{C}^3_\text{col}$.

The framework's terminal prediction:

\begin{itemize}
\item The lattice's natural symmetry at the discrete tick scale
is the left-right symmetric algebra
$\mathfrak{su}(6)_+ \oplus \mathfrak{su}(6)_- \oplus \mathfrak{u}(1)$
of Eq.~\eqref{eq:su6_su6_u1}, with the bipartite parity acting as
the $L \leftrightarrow R$ exchange.  This is the algebraic shape
of a left-right symmetric grand unification, derived from the
$\mathcal{A}=1$ constraint on the bipartite octahedral lattice
rather than postulated.

\item The SM gauge group $SU(3) \times SU(2) \times U(1)$ is the
factor-product subalgebra of this larger structure, recovered
under the (interpretive) restriction that gauge invariance respects
the tensor decomposition.  The 59 ``extras''
(Eq.~\eqref{eq:extras_decomp}) are coset generators that
transform as $(\mathbf{8}, \mathbf{3})$ leptoquark-flavoured plus
chirality-shadow SM under the SM adjoint action.

\item The lattice's bipartite parity is intrinsically linear
(\emph{spatial} parity), and the proposed $SU(2)_W$ couples
\emph{vector-like} to matter, not chirally.  The SM's chirality
and $CP$ violation are not derivable from the framework at the
discrete level --- they are external choices that the SM's
Lagrangian formalism imposes by separating kinetic and mass terms.
The framework's mass mechanism (mass-as-clock-density,
\nolinkurl{notes/mass_chirality_coupling.md}) is structurally
incompatible with chiral gauge coupling.

\item Gravity is a scalar clock-density field on the flat
bipartite lattice (Paper~I~§7), not a curved spacetime metric.
Spacetime torsion --- the antisymmetric part of the affine
connection in Einstein-Cartan theory --- is therefore not a
separate ingredient the framework requires; its conventional
motivation (coupling Dirac spin to curved gravity) is moot
because the framework provides the Dirac structure intrinsically
through the bipartite tick rule's chirality oscillation, without
ever invoking a spin connection (\nolinkurl{notes/no_spacetime_torsion.md}).
\end{itemize}

\subsection{Quantitative predictions}
\label{subsec:quantitative_predictions}

Three quantitative predictions emerge from the Phase~1-4 work:

\paragraph{Gauge-coupling ratio at the lattice scale.}
Eq.~\eqref{eq:g3_over_g2}: $g_3^2 / g_2^2 = 3/2$ at $1/a$.  The SM
measured ratio at $M_Z$ is $g_3^2/g_2^2 \approx 3.3$; the
factor-of-$\sim 2$ shift is consistent with RG running over many
energy decades (asymptotic freedom of $g_3$, slow running of
$g_2$).  The framework does not yet fix the lattice scale $1/a$ from
first principles, so the prediction is a ratio at an unknown UV
scale; a follow-on paper would extract $1/a$ either from the
explicit one-loop prefactor calculation or from another framework
observable.

\paragraph{Anisotropic photon dispersion along $(1, 1, -1)$ (Paper~I); structural extension to non-abelian.}
Paper~I~§6 / \nolinkurl{induced_gauge_action.tex} established
gauge-sector birefringence aligned with the optical axis
$\mathbf{V}_1 + \mathbf{V}_2 + \mathbf{V}_3 = (1, 1, -1)$ for the
abelian photon, derived from the bipartite-plaquette Q-tensor
with eigenvalues $\{4, 4, 16\}$.  Phase~4's non-abelian
generalisation inherits the same Q-tensor unchanged, so the
non-abelian gauge bosons ($W^\pm$, $Z$, gluons) inherit the same
direction-dependent kinetic structure as a structural corollary
of the bipartite geometry.  We flag this as a structural
consequence, not as an independent quantitative prediction
verified in this paper.

\paragraph{Anisotropic decoherence in the colour-weak sector (hypothesis).}
A structural hypothesis (\nolinkurl{notes/extras_as_a1_accounting.md},
\texttt{DRAFT}): the 59 extras of Eq.~\eqref{eq:extras_decomp} are
the algebraic bookkeeping channels through which $\mathcal{A}=1$
tracks probability during what an SM observer identifies as
decoherence.  Decoherence, in this reading, is the algebraic
projection from the full 71-dim centralizer onto the SM-gauge-
invariant 12-dim subalgebra; the ``lost coherence'' resides in
the 59 extras directions.  The
$(\mathbf{8}, \mathbf{3})$ representation content of the
leptoquark-flavoured extras predicts colour-octet weak-triplet
anisotropy in decoherence rates --- a specific testable departure
from continuum decoherence theory.  Verification requires the
quantitative follow-on calculation sketched in
\nolinkurl{notes/extras_as_a1_accounting.md}; this paper does not
execute it.

\subsection{Methodological reflections}
\label{subsec:methodological}

The framework's contribution is sharper than ``a derivation of the
SM gauge group from a discrete substrate.''  Paper~II identifies
\emph{precisely which} features of the Standard Model are geometric
consequences of $\mathcal{A}=1$ on $\mathcal{T}_\diamond^3$ and
\emph{which are not}.

\paragraph{Geometric:} the Lie algebra structure
$\mathfrak{so}(3,1) \oplus \mathfrak{su}(3) \oplus \mathfrak{su}(2)
\oplus \mathfrak{u}(1)$ (dim 18) embedded in a larger left-right
symmetric algebra (dim 71); the bipartite plaquette gauge
invariance and the Q-tensor with eigenvalues $\{4, 4, 16\}$; the
discrete spatial $O_h$ point group (order 48); the colour
$\mathfrak{su}(3)$ generated dynamically from the
$\mathbb{C}^3$ memory.

\paragraph{Not geometric (external SM choices):} chirality (the SM
selects $\sigma_z$ as the chirality axis where the lattice has
$\sigma_x$ as bipartite parity --- orthogonal Bloch involutions);
$CP$ violation (no antilinear modification of the lattice tick
rule admits a Branch~B $CP$); the Higgs mechanism separating
kinetic and mass terms (the lattice's mass \emph{is} the kinetic
mechanism, via the chirality-mixing $\cos(\delta\phi/2)$
amplitude); the factor-product gauge invariance (the lattice
admits 59 factor-mixing operators that the SM forbids by
construction).

\paragraph{The substrate-first approach.}  The framework begins
from a single conservation axiom on a discrete substrate and
works out what physics follows.  This is different from
constructing a model around observed phenomena and fitting
parameters; the framework's predictions are derived from the
lattice geometry without post-hoc parameter fitting.  Paper~I
(\href{https://doi.org/10.5281/zenodo.20078529}{DOI:
10.5281/zenodo.20078529}; see its audit table for the PASS rows)
demonstrated this by deriving the Dirac equation, the Born rule,
hydrogen quantization at four significant figures, and the
structural form of the induced gauge action.  Paper~II continues
the programme, identifying the Lie-algebra-level features that
the substrate produces and the Lagrangian-level features it does
not.

\paragraph{The debt to measurement.}  That the framework's
derivations land where measurement has spent a century pinning
the answers down is not a coincidence --- it is the prerequisite.
A century of accurate, well-tested physics is what makes
substrate-first derivation feasible in the first place: there
must be something firm to land on.  The framework's contribution
is not to replace conventional physics but to identify which of
its features are geometric consequences of a single conservation
law and which are not.  Both readings honour what the testing
community established; the framework reads the established answers
as emergent rather than postulated, but it has no answers to
offer about features the established answers have not already
pinned down (\nolinkurl{notes/debt_to_measurement.md}).

\paragraph{The audit-table convention.}  The audit-table discipline
inherited from Paper~I (Section~\ref{subsec:audit_convention}) is
what makes documentation of emergence credible rather than
aspirational.  Every claim of the paper traces to an artifact:
a sympy-verified statement in \nolinkurl{src/utilities/}, a
derivation in the text, or a measured value cited from the
literature.  Drift between drafts is mechanically detectable;
negative results stay visible (Phase~1.5's \texttt{FAIL} row on
SM-style $CP$ is part of the paper's contribution, not an
embarrassment to be elided); the canonical record is the table,
not the prose.  The convention is a small mechanical change with a
substantial accountability shift, and we hope it is useful to
other first-principles theoretical proposals.

\subsection{Open directions and follow-on papers}
\label{subsec:open_directions}

Several follow-on calculations are scoped but not executed in
this paper.

\paragraph{The universal one-loop prefactor $c$.}  Paper~I left
the explicit $-\mathrm{Tr}\,\ln D_\text{lat}[U]$ calculation as a
``paper-length follow-on''; Paper~II inherits this open item.
Closing it would fix the absolute scale of $1/g^2(1/a)$, enabling a
quantitative RG-running check against the measured
$g_i^2(M_Z)$ values.

\paragraph{The lattice scale $1/a$.}  The framework's lattice-scale
predictions assume a UV cutoff that is not yet fixed from first
principles.  Possible inputs: the Planck scale, the hydrogen
quantization scale, or the framework's induced gravitational
coupling.  Fixing $1/a$ converts the ratio prediction
$g_3^2/g_2^2 = 3/2$ into a quantitative test against measured
couplings at any energy scale via RG running.

\paragraph{The decoherence calculation.}  The structural
hypothesis captured in \nolinkurl{extras_as_a1_accounting.md}
converts the 59 extras into a predictive structure for
anisotropic decoherence rates.  A follow-on script
\nolinkurl{decoherence_as_projection.py} (sketched in
the note) would compute the rate of probability flow from
SM-projected to extras-projected wavefunctions as a function of
the tick rule's $\delta\phi$ and the initial coherence,
converting the qualitative hypothesis into a quantitative
prediction.

\paragraph{The symmetry-breaking mechanism.}  The 71-dim
centralizer at the lattice scale reduces to the SM's 12-dim
gauge group at observable energies.  By analogy with standard GUT
breaking, the 59 extras would correspond to massive gauge bosons
that decouple at lower energies.  Identifying the framework's
breaking mechanism --- whether it is Higgs-analogous, decoherence-
driven via the algebraic projection of
\nolinkurl{notes/extras_as_a1_accounting.md}, or geometric in
some other sense --- is the natural next paper in the series.

\paragraph{Follow-on papers in the catalogue.}  The downstream
papers indexed in
\nolinkurl{external/dcl/notes/follow_on_implications.md} that depend
on Paper~II's colour machinery include item~\#16's proton-
interior study, item~\#3's Farey-sequence mass spectrum, and
items \#13 (Operation Algebra of the DCL) and \#14 (Balanced
Equations and Birefringent Channels) on which the framework's
notational and algebraic infrastructure rests.  Paper~II's
results enter those papers as imported infrastructure: the SU(3)
representation content of colour-memory amplitudes, the
left-right symmetric extension, and the quantitative
gauge-coupling ratios at the lattice scale.

\subsection{Closing}
\label{subsec:closing}

Eq.~\eqref{eq:automorphism_conjecture} of Paper~I asked whether the Standard Model gauge group
plus Lorentz emerges from a single conservation axiom on the
bipartite octahedral causal lattice.  The answer, after four
phases of work: yes as containment, with a precise structural
characterisation of the gap to equality.  The framework realises
a left-right symmetric algebra $\mathfrak{su}(6) \oplus
\mathfrak{su}(6) \oplus \mathfrak{u}(1)$ at the lattice scale,
with the SM gauge group as its factor-product projection.  The
framework's quantitative prediction $g_3^2/g_2^2 = 3/2$ at the
lattice scale, the bipartite-plaquette birefringence aligned with
the $(1, 1, -1)$ optical axis, and the anisotropic decoherence
hypothesis are the testable consequences this work hands to
experiment.

What the framework derives is the Lie-algebra structure of the SM
gauge group from substrate geometry; what it does not derive is
the SM's chirality, $CP$ violation, or Higgs mechanism.  Both
readings are concrete contributions: the geometric features are
recovered without parameter fitting, and the non-geometric
features are identified as external choices in the SM's
Lagrangian formulation that the substrate does not impose.  The
discrete causal lattice with $\mathcal{A}=1$ does not reach the
full Standard Model, but it reaches the maximal structure a
substrate-first derivation can produce, and it identifies the
external inputs the SM requires beyond the substrate.
